\documentclass{EmergentUDB}
\usepackage[utf8]{inputenc}
\usepackage[T1]{fontenc}
\usepackage{geometry}
\usepackage{graphicx}
\usepackage{float}
\usepackage{hyperref}
\usepackage[english,indonesian]{babel}

\title{Optimasi Pemrosesan Template LaTeX Menggunakan Pendekatan Berbasis AI}
\author{Academic Support Agent}
\date{\today}

\begin{document}
\maketitle

\begin{abstract}

\end{abstract}

\chapter{PENDAHULUAN}


\chapter{TINJAUAN PUSTAKA}


% Original content (commented out):
% %%%%%%%%%%%%%%%%%%%%%%%%%%%%%%%%%%%%%%%%%%%%%%%%%%%%%%%
% Official Template of the Emergent Journal
% Universidad Don Bosco, El Salvador
% Please note that this template shows a preview of the article and is not necessarily the final appearance
% For any inquiries, please contact: Manuel Cardona, manuel.cardona@udb.edu.sv



%-------------------------------------------
\documentclass[letterpaper]{EmergentUDB}
%-------------------------------------------
\usepackage{graphicx}
\usepackage{siunitx}
\usepackage{amsmath,amssymb,amsfonts}
\usepackage{listings}
\usepackage{graphicx}
\usepackage{textcomp}
\usepackage{url}
\usepackage{comment}
\def\UrlFont{\rmfamily}
\usepackage{booktabs}
\usepackage{multicol}
\usepackage{float} 
\usepackage{multirow, array}
\usepackage[utf8]{inputenc} 
\usepackage[sorting=none]{biblatex}
\usepackage{academicons}
\addbibresource{biblio.bib}
\usepackage{svg}
\usepackage{orcidlink}


\journal{jcudb}
\title{Emergent Journal Article Template}
\Enlace{http://hdl.handle.net/11715/960}
\jyear{202X}
\jvolume{X}
\jnumber{X}
\jstartpage{X}
\jendpage{X}
\jpubdate{Day month year}


\author[1{\authfn{1}}]{First Author \orcidlink{0000-0000-0000-0000}}
\author[2]{Second Author \orcidlink{0000-0000-0000-0000}}
\author[3]{Third Author\orcidlink{0000-0000-0000-0000}}

\affil[1]{First Institution}
\affil[2]{Second Institution}
\affil[3]{Third Institution}

%%% Autor para Correspondencia
\authnote{\authfn{1}Corresponding author. E-mail: \href{mailto: abc@uni.edu}{emergent.journal@udb.edu.sv, manuel.cardona@udb.edu.sv}}

%%% Categoría del paper
\papercat{Research Article}
%Editorial
%Research Article
%Review Article

\runningtitle{Emergent Journal, Universidad Don Bosco}


\begin{document}
\begin{frontmatter}{a}{b}{c}
\maketitle
\begin{abstract}
Paper should be written in English or Spanish. However, if the the article is written in Spanish, Title, Abstract and Keywords must be written in English. The abstract should frequently indicate the purpose of the study, main results and main conclusions. An abstract is often submitted separately from the manuscript, so it should stand alone. For this reason, references should be avoided. Also, non-standard or unusual abbreviations should be avoided. 
\end{abstract}

\begin{keywords}
%Incluya hasta 6 palabras clave en Inglés.
Keyword 1, keyword 2, keyword 3, keyword 4, keyword 5, keyword 6 (Up to 6 keywords, written in English)
\end{keywords}
\end{frontmatter}

\section{Introducción}

This is the \LaTeX{} template for submitting scientific articles to the \textbf{Emergent} journal of Universidad Don Bosco, El Salvador. Please note that while this template provides a preview of the manuscript for submission, it may not necessarily be the final publication layout.

There are important commands in the preamble that you will need to modify for your own manuscript. If you are using this template in Overleaf, switch the editor to source code mode to view them; or if you prefer to stay in rich text view, click on the title in rich text view to display the preamble code.

Please select the type of article (Research Article or Review Article) in the preamble, with the command"\verb|\papercat{...}|.


\section{Journal Policies}
\subsection{Authenticity}
All articles submitted to \textbf{Emergent} must be original and not be under consideration for publication elsewhere. In case evidence of similar or identical previously published documents is detected, the manuscript will be automatically rejected. Manuscripts with a plagiarism matching percentage greater than or equal to 20\% are not candidates for publication.

\subsection{Peer Review Process}
\textbf{Emergent} journal uses the single-blind peer review system. The editor conducts an initial assessment of the received articles to verify their authenticity, compliance with the journal's objectives, and format. If the article meets these criteria, the editor sends it to three independent reviewers who evaluate its scientific quality.

The editor is responsible for the final decision on whether to accept the article for publication or not. The journal publishes scientific-technical articles, communications or short articles, monographs, and review articles in all areas of engineering and technology. For an article to be accepted, it must represent a significant advancement in the state of the art and demonstrate a solid scientific or technical foundation.

\subsection{Publication Frequency}
\textbf{Emergent} journal publishes biannually (1 volume per year with 2 issues: July and December).

\subsection{Open Access Policy}
\textbf{Emergent} is an Open Access journal; all published articles can be freely downloaded and are available under the Creative Commons license Creative Commons Attribution-NonCommercial-NoDerivatives (CC BY-NC-ND). For non-commercial purposes, it allows others to distribute and copy the article. It also allows its inclusion in a collective work, provided that the original work is cited and the article is not altered or modified.

\subsection{Article Length}
Articles should have a maximum length of 6 pages (including references and annexes) for research articles or 12 pages for review articles (state-of-the-art compilation). The minimum length for a paper is 4 pages.


\section{Submission Process}
Articles must be submitted through the platform. To do so, authors must register as authors (if it is the first time) or log in and access the author profile (if already registered). From there, they can start the submission process. Submissions must be in PDF format.

The document must include a minimum of 250-300 words abstract and must be submitted using the journal template available in the website of the journal. Articles will be reviewed by members of the editorial board based on technical quality, originality, significance, clarity, and scope of research.

The decision on acceptance, revision, or rejection of the manuscript will be communicated solely to the corresponding authors' email.

The journal welcomes contributions with research in the areas of engineering and technology from different disciplinary perspectives. Journal topics include, but are not limited to:

```latex
\begin{itemize}
    \item Embedded Systems
    \item Biomedical
    \item Sensors
    \item Robotics
    \item Bioinformatics
    \item Artificial Intelligence
    \item Natural Language Processing
    \item Machine Learning
    \item Deep Learning
    \item Cyber-Physical Systems
    \item Internet of Things
    \item Communication Networks
    \item Intelligent Manufacturing
    \item Cloud Computing
    \item Digital Transformation
    \item Smart Cities
    \item Industry 4.0
    \item Disruptive Technologies
    \item Virtual Reality
    \item Augmented Reality
    \item Metaverse
    \item Chatbots
    \item Cybersecurity
    \item Unmanned Aerial Vehicles
    \item Computer Vision
    \item STEAM Initiatives
\end{itemize}


\subsection{Submission Checklist}

Make sure to comply with the following guidelines before submitting your work:

\begin{itemize}
    \item[\checkmark] As part of the submission process, authors are required to ensure that their submission meets all the elements shown below. Works that do not comply with these guidelines will be returned to the authors.
    
    \item[\checkmark] The journal's policy requires authors to submit only original articles that have not been previously published or be under review for another peer-reviewed publication. The manuscript has a plagiarism matching percentage of less than 20\%.
    
    \item[\checkmark] Articles can be submitted in English or Spanish. If submitted in Spanish, the title, abstract, and keywords must be written in English.
    
    \item[\checkmark] The maximum length of the research article should be 6 pages (including references and annexes), and 12 pages for review articles. The minimum number of pages should be 4.
    
    \item[\checkmark] The article must strictly adhere to the style and bibliography type described in the template. Template are available in LaTeX and can be downloaded from the following link:
    
    \begin{itemize}
        \item LaTeX template
    \end{itemize}
    
    \item[\checkmark] The article file must be submitted in PDF format, generated from the journal's templates.
    
    \item[\checkmark] The article should contain: Title (in English), Abstract (in English), Keywords (in English), Introduction (in English or Spanish), Methods (in English or Spanish), Results (in English or Spanish), Discussions (in English or Spanish), Conclusions (in English or Spanish), References (in English).
    
    \item[\checkmark] It must include a PNG color image as a graphical abstract for the article. The image should reflect the main contribution of your article as it constitutes a visual summary of the main findings. Provide an image with a minimum of 531 pixels x 1328 pixels (height x width) with a minimum resolution of 300 dpi.
\end{itemize}

\section{Article Format Considerations}

\subsection{Figures and Tables}
Figures and tables can be added using the \verb|figure| and \verb|table| environments, for example, Fig. \ref{fig:ejemplo} and Table \ref{tabla:ejemplo}.


\begin{figure}[!h]
\centering
\includegraphics[width=2.5in]{ALICE.png}
\caption{Exoskeleton Robotic Rehabilitation ALICE \cite{CardonaSensors}}
\label{fig:ejemplo}
\end{figure}

\begin{table}[h]
\centering
\caption{Localization of ALICE's Screws Axes \cite{ScrewsALICE}}
\label{tabla:ejemplo}
\begin{tabular}{ccc}
\toprule
Joint & Axis Direction ($\$_i$) & Location ($s_{0i}$) \\
\cmidrule{1-3}
$1$ & $[0, 0, 1]$ & $[0, 0, 0]$ \\
$2$ & $[1, 0, 0]$ & $[0, 0, l_2]$ \\
$3$ & $[1, 0, 0]$ & $[0, l_3, l_2]$ \\
\bottomrule
\end{tabular}
\end{table}

Use \verb|figure*| and \verb|table*| if you need a two-column wide figure or table, as in Fig. \ref{fig:ejemplowide} and Table \ref{tabla:ejemplowide}.

\begin{figure*}[!h]
\centering
\includegraphics[width=6.4in]{FigSimulink.eps}
\caption{Equivalent Model in Simscape Multibody\textsuperscript{TM} \cite{CardonaSensors}}
\label{fig:ejemplowide}
\end{figure*}

\begin{table*}[!t]
\caption{Selection of ALICE's Actuation System \cite{CardonaDyn}}
\begin{center}
\begin{tabular}{ccccccccc}
\hline
\label{tabla:ejemplowide}
\textbf{Joint} & \textbf{Motor} & \textbf{Reducer} & \textbf{Controller} & \textbf{Voltage (V)} & \textbf{Torque (N.m)} & \textbf{Peak Torque (N.m)} & \textbf{Reducer} \\
\hline 
J1 & EC Flat 90 & HD CSG 14 2UH & EPOS2 70/10 & 24 & 14 & 58 & 80 \\
J2 & EC Flat 90 & HD CSG 17 2UH & EPOS2 70/10 & 24 & 51 & 109 & 100 \\
J3 & EC Flat 90 & HD CSG 17 2UH & EPOS2 70/10 & 24 & 51 & 109 & 100\\                       
\hline 
\end{tabular}
\end{center}
\end{table*}


Remember to insert figures and tables after citing them in the text. Use "Fig. 1", not "Figure 1" or "Figura 1".

For figure labels: Use Times New Roman 8 points. Use words instead of symbols or abbreviations when writing labels for the axes of graphs to avoid confusing the reader. For example, write the quantity "Acceleration", not just "a". If you include units in the label, present them in parentheses. Do not label the axes only with the units. In the example, write "Acceleration ($m/s^2$")".

\subsection{Equations}
Equations must be numbered consecutively, and the equation number should be in parentheses. The equation should be centered, and the numbering aligned to the right. Ensure that the symbols in the equation are defined before or immediately after the equation. Use "(1)", not "Eq. (1)", "Equation (1)", "Eq. 1", or "equation (1)", except at the beginning of a sentence: "Equation (1) is..".

For example: The inverse Fourier transform is given by (\ref{eq:invFourier}).

\begin{equation} \label{eq:invFourier}
    F[X(\omega)]^{-1}=x(t)=\frac{1}{2\pi}\int_{-\infty}^{\infty}X(\omega) e^{j\omega t} d\omega
\end{equation}

While (\ref{eq:matrizskew}) shows the antisymmetric matrix

\begin{equation} \label{eq:matrizskew}
\tilde{s}=\begin{bmatrix}0&-s_z&s_y\\ s_z&0&-s_x\\-s_y&s_x&0\end{bmatrix}
\end{equation}

\subsection{Text Style}
\begin{epigraph}{Source Name of the Epigraph}
This is the text of an epigraph, should you like to add one. Lorem ipsum dolor sit amet, consectetur adipiscing elit, sed do eiusmod tempor incididunt ut labore et dolore magna aliqua. 
\end{epigraph}

\begin{quote}
This is a quote. Lorem ipsum dolor sit amet, consectetur adipiscing elit, sed do eiusmod tempor incididunt ut labore et dolore magna aliqua. Ut enim ad minim veniam, quis nostrud exercitation ullamco laboris nisi ut aliquip ex ea commodo consequat. 
\end{quote}

\begin{itemize}
\item This is a bullet list.
\item Second point.
\item Third point.
\end{itemize}

\begin{itemize}
\item [\checkmark] This is a bullet list with a special character.
\item [\checkmark] Second point.
\item [\checkmark] Third point.
\end{itemize}

\begin{enumerate}
\item This is a numbered list.
\item Second item.
\item Third item.
\end{enumerate}

\begin{mdframed}[linewidth=1pt]
This is a text box. Lorem ipsum dolor sit amet, consectetur adipiscing elit, sed do eiusmod tempor incididunt ut labore et dolore magna aliqua.
\end{mdframed}


\subsection{Citations and References}
Make sure that all references cited in the text are also present in the reference list (and vice versa). Citations in the text should follow the referencing style used by the Institute of Electrical and Electronics Engineers (IEEE). Use the file \verb|biblio.bib| to input the references.

\subsubsection{Examples of References}
Consider these style examples for citing references. Use the \verb|\cite| command to add a citation. For example, this is a citation: \cite{Jackson}, while these are two citations in a single command: \cite{Grady,Piumsomboon:2017}.
\begin{itemize}
\item Reference to a journal: \cite{Grady}
\item Reference to a journal with manuscript number: \cite{Jerrentrup}
\item Reference to a book: \cite{Jackson}
\item Reference to a chapter in an edited book: \cite{Aron, Cardona2020}
\item Reference to a website: \cite{testWeb}
\item Reference to a dataset: \cite{testDataset}
\item Reference to conference proceedings: \cite{Piumsomboon:2017}
\item Reports: \cite{testReport}
\item Reference to dissertations and theses: \cite{testThesis}
\end{itemize}


\subsection{Example of Mathematical Notation}

Let $X_1, X_2, \ldots, X_n$ denote a sequence of independent random variables with $\text{E}[X_i] = \mu$ and $\text{Var}[X_i] = \sigma^2 < \infty$. Then:
%
\begin{equation}
S_n = \frac{X_1 + X_2 + \cdots + X_n}{n}
      = \frac{1}{n}\sum_{i}^{n} X_i
\end{equation}
%
represents the average. Then, as $n$ tends to infinity, the random variables $\sqrt{n}(S_n - \mu)$ converge to $\mathcal{N}(0, \sigma^2)$.

\subsection{Algorithms, Program Codes, and Listings}\label{sec7}

The packages \verb+algorithm+, \verb+algorithmicx+, and \verb+algpseudocode+ are used for writing algorithms in \LaTeX.
For this purpose, the following format should be used:

\begin{verbatim}
\begin{algorithm}
\caption{<alg-caption>}\label{<alg-label>}
\begin{algorithmic}[1]
. . .
\end{algorithmic}
\end{algorithm}
\end{verbatim}

For example, a fast exponentiation procedure is:
 
\begin{algorithm}[!h]
\caption{Calculate $y = x^n$}\label{algo1}
\begin{algorithmic}[1]
\Require $n \geq 0 \vee x \neq 0$
\Ensure $y = x^n$
\State $y \Leftarrow 1$
\If{$n < 0$}
        \State $X \Leftarrow 1 / x$
        \State $N \Leftarrow -n$
\Else
        \State $X \Leftarrow x$
        \State $N \Leftarrow n$
\EndIf
\While{$N \neq 0$}
        \If{$N$ is even}
            \State $X \Leftarrow X \times X$
            \State $N \Leftarrow N / 2$
        \Else[$N$ is odd]
            \State $y \Leftarrow y \times X$
            \State $N \Leftarrow N - 1$
        \EndIf
\EndWhile
\end{algorithmic}
\end{algorithm}

For program codes, the \verb+program+ package should be used; moreover, it's necessary to employ the \verb+\begin{program}+ \verb+...+ \verb+\end{program}+ environment to write the program code.

Similarly, for program codes, the \verb+listings+ package can be used. Additionally, \verb+\begin{lstlisting}+ \verb+...+ \verb+\end{lstlisting}+ is used to establish the environment, similarly for the \verb+verbatim+ environment. For more details, refer to the documentation of the \verb+lstlisting+ package. For example, below is a snippet of code in Python.

\begin{lstlisting}[language=python]
import pandas as pd
import seaborn as sns
import matplotlib.pyplot as plt
import tensorflow as tf
\end{lstlisting}


\section{Article Structure}
While there is no standard format for the sections that an article should contain, it is suggested to consider the elements listed below. Note that each section should be specified with \verb|\section{}|, except for the subsections (as they should not be numbered): Appendices, Acknowledgments, Declaration of Conflicts of Interest, and References; in this case, use the command \verb|\section*{}|.

\section*{-Introduction}
This section sets out the objectives of the work and provides background, avoiding a detailed review of the literature or a summary of the results.

\part{-Materials and Methods}
This section provides sufficient details to allow an independent researcher to reproduce the work. Methods that are already published should be summarized and indicated by a reference. If directly quoting a previously published method, use quotation marks and also cite the source. Any modifications to existing methods should also be described.

The methods section should include the study design, the type of materials involved, a clear description of all comparisons, and the type of analysis used, to allow for the replication of the work. Ease of reproducibility is one of the key criteria upon which reviewers will be asked to comment.

\part{-Results}
This section should comprise the results of the study presented in a logical sequence, complemented by tables and/or figures. Ensure that the text does not repeat data presented in the tables and/or figures.

\part{-Discussion}
This section should explore the meaning of the results of the work, highlighting why they are important. A combined Results and Discussion section is often appropriate. Avoid extensive citations and discussion of published literature.

\part{-Conclusions}
The main conclusions of the study can be presented in a brief Conclusions section, which can be independent or form a subsection of a Discussion or Results and Discussion section.

\part{-Additional Sections}
Authors can add or remove section titles according to their manuscript's type or structure, such as the Appendices section.

\part{-Appendices}
Although optional, an Appendices section can be useful for expanding explanations of procedures, equations, tables, figures, etc.

\part{Acknowledgments}
Write your acknowledgments to funding or support organizations.

\part{Declaration of Conflicts of Interest}
Authors will declare that there is no potential conflict of interest regarding the research, authorship, and/or publication of this article.



\printbibliography

%If the article is written in Spanish please use:
%\printbibliography[title={Referencias}]


\end{document}

\end{document}